% !TeX program = lualatex
% ================================================================
% Polish Parallel Text Version of FractionsDecimalsAndPercentagesConversions
%
% Generated by the server using the following settings:
%
%   language            Polish (pol)
%   text direction      left to right
%   font                Noto Serif
%   fallback font       Noto Serif
%   built from          latex/worksheets/percentages/fractionsDecimalsAndPercentagesConversions.tex
%   vocabulary key      latex/worksheets/percentages/fractionsDecimalsAndPercentagesConversions/L2/pol/fractionsDecimalsAndPercentagesConversions_polishKey.tex
%   tier 3 vocabulary   green   RGB 0 128 0
%   English text        black   RGB 0 0 0
%   Polish text         purple  RGB 112 48 160
%   layout macros       version 3
%
% Please don't edit any information in this comment block - the server
% compares it against the current settings and rebuilds this file when
% they differ, so changes here would be overwritten. However, feel free
% to make updates to the LaTeX below if you want to edit it.
% ================================================================

\documentclass[12pt]{article}
% ================================================================
% Copied from the English sheet this was translated from, so that
% anything it sets up for itself still works here
% ================================================================
\usepackage[a4paper,margin=1in]{geometry}
\usepackage{amsmath}
\usepackage{amssymb}
\usepackage{siunitx}
\usepackage{enumitem}
\usepackage{multicol}

% Characters this language's font has not got are borrowed from another,
% rather than being silently dropped - see docs/translated-sheet-latex.md
\directlua{%
  if luaotfload and luaotfload.add_fallback then
    luaotfload.add_fallback("eallatinfallback",
      { "Noto Serif:mode=harf;" })
  else
    tex.error("this luaotfload is too old to register a fallback font")
  end
}
\usepackage[english]{babel}
\babelprovide[import]{polish}
\babelfont[polish]{rm}[Renderer=Harfbuzz, BoldFont={Noto Serif}, ItalicFont={Noto Serif}, BoldItalicFont={Noto Serif}, RawFeature={fallback=eallatinfallback}]{Noto Serif}
\babelfont[polish]{sf}[Renderer=Harfbuzz, BoldFont={Noto Serif}, ItalicFont={Noto Serif}, BoldItalicFont={Noto Serif}, RawFeature={fallback=eallatinfallback}]{Noto Serif}
\newcommand{\ealtext}[1]{\foreignlanguage{polish}{#1}}
\newcommand{\ealtextblock}[1]{{\raggedright\foreignlanguage{polish}{#1}\par}}
% ================================================================
% EAL layout helpers
% ================================================================
\usepackage{xcolor}

\definecolor{ealkeycolour}{RGB}{0,128,0}
\definecolor{ealenglish}{RGB}{0,0,0}
\definecolor{ealtwocolour}{RGB}{112,48,160}

\newsavebox{\ealboxen}
\newsavebox{\ealboxtr}
\newlength{\ealglosswd}
\newlength{\ealglossraise}
\setlength{\ealglossraise}{1.15em}

% a tier 3 word, in English and in translation alike
\newcommand{\ealkey}[1]{\textcolor{ealkeycolour}{#1}}
\newcommand{\ealkeytr}[1]{\textcolor{ealkeycolour}{\ealtext{#1}}}

% One translated word above one English word. Built from kernel
% primitives, never a tabular, which would break wherever the sheet
% loads array, booktabs or tabularx. The box is as wide as the wider of
% the two, so a long translation pushes its neighbours aside instead of
% overlapping them, and it claims its full height so a table row or a
% TikZ node makes room for it.
\newcommand{\ealgloss}[2]{%
  \sbox{\ealboxen}{\ealkey{#1}}%
  \sbox{\ealboxtr}{\small\textcolor{ealtwocolour}{\ealtext{#2}}}%
  \setlength{\ealglosswd}{\wd\ealboxen}%
  \ifdim\wd\ealboxtr>\ealglosswd\setlength{\ealglosswd}{\wd\ealboxtr}\fi
  \leavevmode
  \makebox[\ealglosswd][c]{%
    \makebox[0pt][c]{\usebox{\ealboxen}}%
    \makebox[0pt][c]{\raisebox{\ealglossraise}{\usebox{\ealboxtr}}}%
  }%
}

% opens up line spacing around a block of glosses, so the raised
% translations sit clear of the line above
\newenvironment{ealglossed}{\par\linespread{2.1}\selectfont\ignorespaces}{\par}

% a whole translated sentence above its English counterpart. block level,
% so it does not belong inside a tabular cell
\newcommand{\ealpara}[2]{%
  \par\addvspace{0.5\baselineskip}%
  {\color{ealtwocolour}\ealtextblock{#1}}%
  \penalty200
  {\color{ealenglish}#2\par}%
  \addvspace{0.5\baselineskip}%
}

\begin{document}
\pagenumbering{gobble}

\title{}
\date{}
\maketitle

\ealpara{Przeliczanie: \ealkeytr{ułamki dziesiętne}, \ealkeytr{procenty} i \ealkeytr{ułamki}}{Conversions: \ealkey{Decimals}, \ealkey{Percentages}, and \ealkey{Fractions}}

\begin{enumerate}[label=\textbf{Q\arabic*.}]

\item
\ealpara{Zapisz następujące \ealkeytr{ułamki dziesiętne} jako \ealkeytr{ułamek} w \ealkeytr{najprostszej postaci} i jako \ealkeytr{procent}:}{Write the following \ealkey{decimals} as a \ealkey{fraction} in \ealkey{simplest form} and as a \ealkey{percentage}:}
\begin{multicols}{2}
\begin{enumerate}[label=(\alph*)]
    \item \( 0.125 \)
    \item \( 0.203125 \)
\end{enumerate}
\end{multicols}


\item
\ealpara{Zamień następujące \ealkeytr{ułamki} na \ealkeytr{ułamki dziesiętne} i \ealkeytr{procenty}:}{Convert the following \ealkey{fractions} into \ealkey{decimals} and \ealkey{percentages}:}
\begin{multicols}{2}
\begin{enumerate}[label=(\alph*)]
    \item \( \frac{7}{20} \)
    \item \( \frac{33}{320} \)
\end{enumerate}
\end{multicols}

\item
\ealpara{Wyraź następujące \ealkeytr{procenty} jako \ealkeytr{ułamki dziesiętne} i \ealkeytr{ułamki} w \ealkeytr{najprostszej postaci}:}{Express the following \ealkey{percentages} as \ealkey{decimals} and \ealkey{fractions} in \ealkey{simplest form}:}
\begin{multicols}{2}
\begin{enumerate}[label=(\alph*)]
    \item \( 17.5\% \)
    \item \( 57.8125\% \)
\end{enumerate}
\end{multicols}

\end{enumerate}


\bigskip


\begin{enumerate}[label=\textbf{Q\arabic*.}, start=4]

\item

\ealpara{\textbf{\ealkeytr{Punkty bazowe}} to jednostki używane w finansach do opisywania małych zmian \ealkeytr{procentów}.}{\textbf{\ealkey{Basis points}} are units used in finance to describe small changes of \ealkey{percentages}.}

\ealpara{\ealkeytr{Punkt bazowy} to \ealkeytr{procent} z \ealkeytr{procentu}:}{A \ealkey{basis point} is a \ealkey{percent} of a \ealkey{percent}:}
\[
1 \text{ \ealkey{basis point}} = 0.01\%
\]

\ealpara{Często używa się ich do opisywania zmian stóp \ealkeytr{odsetek} lub \ealkeytr{dochodów}.}{They are often used to describe changes in \ealkey{interest} rates or \ealkey{yields}.}

\ealpara{Na przykład wzrost o \( 25 \) \ealkeytr{punktów bazowych} oznacza, że stopa \ealkeytr{odsetek} wzrasta o \( 0.25\% \).}{For example, an increase of \( 25 \) \ealkey{basis points} means the \ealkey{interest} rate goes up by \( 0.25\% \).}

\medskip


\begin{enumerate}[label=(\alph*)]
    \item \ealpara{Zapisz jeden \ealkeytr{punkt bazowy} jako \ealkeytr{ułamek dziesiętny} i \ealkeytr{ułamek} w \ealkeytr{najprostszej postaci}.}{Write one \ealkey{basis point} as a \ealkey{decimal} and a \ealkey{fraction} in \ealkey{simplest form}.}
    \item \ealpara{Zapisz 175 \ealkeytr{punktów bazowych} jako \ealkeytr{ułamek dziesiętny} i \ealkeytr{ułamek} w \ealkeytr{najprostszej postaci}.}{Write 175 \ealkey{basis points} as a \ealkey{decimal} and a \ealkey{fraction} in \ealkey{simplest form}.}
\end{enumerate}

\bigskip

\ealpara{\textbf{\ealkeytr{Brytyjskie obligacje skarbowe}} to \ealkeytr{obligacje} rządowe Wielkiej Brytanii.}{\textbf{\ealkey{Gilts}} are UK government \ealkey{bonds}.}

\ealpara{Kiedy kupujesz \ealkeytr{brytyjską obligację skarbową}, pożyczasz pieniądze rządowi.}{When you buy a \ealkey{gilt}, you are lending money to the government.}

\ealpara{Rząd Wielkiej Brytanii obecnie emituje 30-letnie \ealkeytr{brytyjskie obligacje skarbowe} oferujące prosty \ealkeytr{zwrot} (bez \ealkeytr{kapitalizacji}) w wysokości $5.355\%$ rocznie.}{The UK government currently issues 30-year \ealkey{gilts} offering a simple \ealkey{return} (not \ealkey{compounding}) of $5.355\%$ per year.}

\ealpara{Używają ich do pożyczania pieniędzy na finansowanie inwestycji.}{They use these to borrow money to fund investment.}

\ealpara{Każdego roku płacą \textbf{\ealkeytr{kupon}} równy $5.355\%$ pożyczonej kwoty.}{Each year they pay a \textbf{\ealkey{coupon}} equal to $5.355\%$ of the amount borrowed.}

\ealpara{Po upływie 30 lat spłacają pierwotną pożyczoną kwotę.}{At the end of the 30 years they pay back the original amount borrowed.}

\ealpara{Sprzedawane są również inne okresy (1-letnie, 5-letnie, 10-letnie) \ealkeytr{brytyjskich obligacji skarbowych}.}{Other durations (1-year, 5-year, 10-year) of \ealkey{Gilts} are also sold.}

\medskip

\begin{enumerate}[label=(\alph*), resume]
    \item \ealpara{Znajdź $5.355\%$ jako \ealkeytr{ułamek dziesiętny} i \ealkeytr{ułamek} w \ealkeytr{najprostszej postaci}.}{Find $5.355\%$ as a \ealkey{decimal} and a \ealkey{fraction} in its \ealkey{simplest form}.}
    \item \ealpara{Oblicz całkowity koszt, jaki rząd poniesie, pożyczając £1000 na 30 lat, wraz ze spłatą pierwotnej kwoty.}{Calculate the total cost for the government to borrow £1000 over 30 years, including payback of the original amount.}
    \item \ealpara{Załóż, że płatność \ealkeytr{kuponu} z 30-letniej \ealkeytr{brytyjskiej obligacji skarbowej} wzrasta o \( 43.5 \) \ealkeytr{punktów bazowych}. O ile więcej kosztuje rząd pożyczenie £1000 na 30 lat?}{Suppose the \ealkey{coupon} payment of a 30 year \ealkey{gilt} rises \( 43.5 \) \ealkey{basis points}. How much more does it cost the government to borrow £1000 over 30 years?}
    % \item Gilts offer only simple interest, how might an investor achieve compounding returns, while still only investing in Gilts?
\end{enumerate}

\end{enumerate}

\end{document}